\documentclass[12pt]{article}

\usepackage{sbc-template}
\usepackage[brazil]{babel}
\usepackage{graphicx}
\usepackage{times}
\usepackage{url}

\sloppy

\title{A Metamorfose em vários aspectos: Uma Leitura Pessoal sobre o Adoecimento e a Perda de Si em Franz Kafka}

\author{Yago Melo da Costa}

\begin{document}
\maketitle

\begin{resumo}
Este artigo de opinião apresenta uma análise subjetiva da obra "A Metamorfose", de Franz Kafka, a partir da primeira leitura do autor no ano de 2026. O texto explora a confusão inicial acerca da alienação de Gregor Samsa por seu trabalho, a angústia da perda da comunicação e a reação de sua família diante da sua incapacitação. Através de uma reflexão pessoal, o artigo propõe que a metamorfose física de Gregor é uma poderosa metáfora para o adoecimento mental e físico, a depressão e a consequente perda de utilidade em uma sociedade capitalista. A obra é interpretada como um relato da fragilidade dos laços familiares quando submetidos à pressão financeira e à incapacidade do outro, ecoando experiências pessoais do leitor com o adoecimento de familiares.
\end{resumo}

\begin{abstract}
This opinion article presents a subjective analysis of Franz Kafka's "The Metamorphosis," based on the author's first reading in 2026. The text explores the initial confusion about Gregor Samsa's alienation by his work, the anguish of losing communication, and his family's reaction to his incapacitation. Through a personal reflection, the article proposes that Gregor's physical metamorphosis is a powerful metaphor for mental and physical illness, depression, and the consequent loss of utility in a capitalist society. The work is interpreted as an account of the fragility of family bonds when subjected to financial pressure and the other's incapacity, echoing the reader's personal experiences with the illness of family members.
\end{abstract}

\section{Introdução}

\noindent
\textbf{Edição utilizada:} Todas as citações e referências a números de página neste artigo referem-se à edição bilíngue publicada pelo Instituto Mojo \cite{kafka2022}, disponível gratuitamente em formato digital. Foi a única edição consultada para esta leitura. É importante situar que "A Metamorfose" foi escrita por Franz Kafka e publicada originalmente em 1915 sob o título alemão \emph{Die Verwandlung}. A edição que li é uma tradução para o português realizada por Giovane Rodrigues, publicada em 2022, que mantém o texto bilíngue (português-alemão).

Quando me aventurei pela primeira vez nas páginas de "A Metamorfose" \cite{kafka2022}, não esperava que um conto sobre um homem que vira um inseto me causasse tamanha comoção. A princípio, a obra me pareceu fantástica, quase absurda. No entanto, à medida que acompanhava os monólogos internos de Gregor Samsa e a reação de sua família, fui tomado por uma angústia e uma tristeza profundas. Este artigo é um exercício de opinião, uma tentativa de organizar o que entendi literalmente e, principalmente, o que o texto despertou em mim de maneira subjetiva.

O objetivo não é esgotar as interpretações acadêmicas sobre Kafka, mas sim dialogar com a obra a partir das minhas próprias experiências e sentimentos. Como um leitor que se identifica mais com o cinema do que com a literatura, este foi um livro que me fisgou pela emoção e me forçou a refletir sobre temas como trabalho, adoecimento, família e solidão.

\section{Sobre a Obra e a Primeira Impressão}

Inicialmente, fiquei genuinamente confuso com a reação de Gregor. Ele acorda transformado em um "imenso inseto" \cite[p. 9]{kafka2022}, com um corpo que ele não reconhece e que não obedece aos seus comandos. No entanto, sua maior preocupação não é a transformação em si, mas sim como isso afetará sua capacidade de trabalhar. Ele pensa: "Que profissão cansativa eu escolhi! [...] Que o diabo carregue isso tudo!" \cite[p. 10-11]{kafka2022} e, em seguida, já se preocupa em perder o trem das cinco horas.

Essa atitude me pareceu, num primeiro momento, inacreditável. Como alguém pode se importar mais com um emprego do que com o fato de ter virado um besouro? Aos poucos, fui entendendo que Kafka não estava apenas descrevendo um choque, mas sim o grau máximo de alienação. Gregor estava tão preso à sua rotina de trabalho, motivado pela necessidade de pagar a dívida dos pais com o patrão, que sua própria identidade se dissolveu. Ele não era mais ele mesmo; ele era o "caixeiro-viajante" provedor. A metamorfose apenas tornou visível uma verdade que já existia: ele já era um inseto para o sistema, uma peça descartável e mecânica.

\section{A Angústia do Isolamento}

O segundo ponto que me causou imensa angústia foi a progressiva perda da capacidade de comunicação de Gregor. Ao tentar falar com a mãe e o gerente, ele percebe que sua voz se transformou em um "som esganiçado, dolorido" \cite[p. 13]{kafka2022}. Sua irmã, Grete, é a única que ainda tenta se comunicar, mas logo todos desistem de entendê-lo.

A partir dali, Gregor se torna um observador mudo do próprio destino. Ele tem pensamentos lúcidos, complexos e repletos de sentimentos, mas não há mais uma ponte para expressá-los. Para mim, essa foi a representação mais poderosa da depressão e do adoecimento mental. A sensação de estar preso dentro de si mesmo, de ter tudo a dizer mas nenhuma ferramenta para fazê-lo, e de ver as pessoas ao redor agindo com base em seus próprios medos e nojo, sem nunca realmente saber o que se passa no outro. A cena em que ele tenta, com a boca, girar a chave na fechadura, machucando-se e soltando um "fluido marrom" \cite[p. 28]{kafka2022}, é a imagem perfeita desse esforço inglório e doloroso de tentar, por qualquer meio, ser novamente parte do mundo humano.

\section{A Reação da Família}

A dinâmica familiar, para mim, foi o que tornou a leitura tão cheia de camadas e, ao mesmo tempo, me causou raiva e tristeza. Cada membro reagiu de um jeito, o que é possível explicitar na minha opinião de uma primeira leitura cada comportamento dos membros da família Samsa.

Começo pelo pai, o Sr. Samsa. Ele é um homem moldado para obedecer. Antes da metamorfose, vivia derrotado, falido, dependente do filho, usando um velho roupão e se arrastando com muletas \cite[p. 67]{kafka2022}. A transformação de Gregor o força a sair dessa inércia. Ele arruma um emprego no banco e resgata uma autoridade que julgava perdida, o uniforme que ele se recusa a tirar mesmo em casa é o símbolo dessa nova identidade rígida \cite[p. 72-73]{kafka2022}. No entanto, sua reação a Gregor é bem violenta. O ponto de virada é o ataque com maçãs: ele bombardeia o filho indefeso, e uma maçã fica cravada nas costas de Gregor, causando uma ferida que nunca cicatriza \cite[p. 69-70]{kafka2022}. Mas o que torna o pai tão complexo é que sua violência está ligada à sua própria recuperação. Livrando-se do fardo que Gregor representava, ele perde o medo de não obedecer ou ser devoto a um patrão e, no final, expulsa os hóspedes com firmeza \cite[p. 96-97]{kafka2022}.

A mãe, por sua vez, é a personagem que mais me gera indecisão. Ela demonstra amor, mas sua fragilidade a impede de agir. É ela quem defende que o quarto de Gregor não seja esvaziado, acreditando que isso o desencorajaria \cite[p. 59]{kafka2022}, e ainda pede que poupem a vida do filho durante o ataque do pai \cite[p. 70]{kafka2022}. No entanto, esse amor é interno, silencioso e não exposto totalmente. Ela sofre de asma e desmaia sempre que confrontada com a realidade de Gregor na primeira vez que o vê \cite[p. 29]{kafka2022} e novamente ao vê-lo na parede \cite[p. 64]{kafka2022}. O ápice da ambiguidade está no final: após a morte de Gregor, o pai diz “agora podemos agradecer a Deus” \cite[p. 95]{kafka2022}, e a mãe, em vez de lamentar, sente alívio. 

A irmã Grete é a personagem que mais me causou raiva e decepção, justamente porque sua trajetória é a que teve maior desenvolvimento na família Samsa. No início, ela é a única cuidadora ativa: limpa o quarto \cite[p. 44]{kafka2022}, tenta descobrir o que Gregor gosta de comer e age como sua principal conexão com o mundo. No entanto, à medida que a situação financeira se agrava e ela precisa trabalhar, sua compaixão se esgota. A grande virada ocorre durante o concerto de violino: ao ver Gregor saindo do quarto, ela se sente humilhada na frente dos hóspedes e passa a se referir a ele como “essa coisa” \cite[p. 89]{kafka2022}. É ela quem profere a sentença final: “Precisamos nos livrar disto aí” \cite[p. 89]{kafka2022}. Ao matar simbolicamente o irmão, Grete encontra sua própria força. Ela, que antes era vista como uma adolescente inútil, torna-se uma mulher trabalhadora e com "voz" de agir, visto que, semelhante a Gregor antes da metamorfose, age como elo da família Samsa. No desfecho, é descrita como uma “moça bonita e viçosa” \cite[p. 100]{kafka2022}, algo que antes nunca foi exposto.

Um contraste interessante vem da faxineira, uma figura externa à família. Diferentemente dos Samsa, ela não sente nojo nem medo de Gregor. Desde o primeiro dia em que o vê, chama-o de “besouro-do-esterco” e trata-o como um bicho comum, sem a carga emocional que a família carrega \cite[p. 78]{kafka2022}. Ela não se assusta com a sua presença e chega até a provocá-lo. Esse comportamento, aparentemente rude, evidencia algo crucial: a reação de pavor e rejeição dos Samsa não era inevitável. Era uma escolha. A faxineira, justamente por não ter nenhum vínculo afetivo com Gregor, consegue lidar com ele de forma muito mais natural e menos violenta do que a própria família.

Sinto raiva desses personagens, mas não posso desconsiderar o contexto que a família passava. A situação financeira era terrível e eles estavam desesperados. No entanto, como leitor que tem acesso à mente de Gregor, é impossível não se indignar com a forma como o sofrimento dele foi tratado como um mero incômodo.

\section{Minha Interpretação Subjetiva}

Minha percepção central da obra é que a metamorfose física é um recurso literário para expressar o que acontece quando uma pessoa se torna "incapacitada". Para mim, não importa se Gregor virou um besouro, uma barata ou um escaravelho. O que importa é que ele acordou um dia e não pôde mais cumprir o seu papel social de trabalhador e provedor.

Vejo isso como uma metáfora forte para o adoecimento. Pode ser um coma, uma depressão profunda, uma crise de ansiedade debilitante ou uma doença degenerativa como o Alzheimer. A pessoa ainda está ali, em algum nível, mas não consegue mais se expressar, trabalhar ou ser útil nos moldes da sociedade. A reação da família com sentimentos como o nojo, a raiva, a negligência e, finalmente, o desejo de se "livrar" do problema. Infelizmente, é um comportamento humano mais comum do que se imagina.

\section{Reflexões Pessoais e Conclusão}

Ler "A Metamorfose" foi difícil sendo sincero. Mexeu profundamente comigo. Lembrei imediatamente da minha avó, que ficou incapacitada pelo Alzheimer. Vi a minha própria família paterna tendo que se reconstruir, as brigas, as mudanças de vida e, sim, aquele sentimento velado de que a pessoa já não estava mais ali, sendo apenas um fardo que tinham que cuidar até sua partida. O livro também me fez lembrar de fases da minha própria vida em que a depressão me confinou a um quarto, me tornando um observador apático do mundo, exatamente como Gregor sob o canapé.

A tristeza que senti foi imensa, pois me identifiquei com seus monólogos e com a sensação de habitar um corpo que não parece seu, limitado e incompreendido. Isso demonstra, para o meu contexto pessoal, ser um obra bastante imersiva.

Este foi o meu primeiro livro lido em 2026, e como alguém que sempre preferiu filmes, ele reacendeu em mim o desejo de buscar mais obras que me façam refletir, evoluir e, acima de tudo, sentir. Escolhi um livro com um teor melancólico, sem saber o impacto que teria. Agora, seguirei para uma leitura mais leve para digerir todas as camadas que Kafka me apresentou. Apenas posso concluir que, para mim, o texto fala sobre a fragilidade dos laços humanos quando a materialidade e a produtividade deixam de existir, e sobre a solidão de quem, por qualquer motivo, é excluído do ciclo produtivo.

\begin{thebibliography}{99}

\bibitem{kafka2022}
KAFKA, Franz.
\emph{A Metamorfose} [\emph{Die Verwandlung}, 1915].
Tradução de Giovane Rodrigues.
São Paulo: Instituto Mojo, 2022. (Coleção Literatura Livre).
Edição bilíngue Português-Alemão. E-book.
ISBN 978-65-89008-23-1.
Disponível em: \url{https://mojo.org.br/wp-content/uploads/2024/08/006-A-metamorfose-978-65-89008-23-1.pdf}.

\end{thebibliography}

\end{document}